% Historical synthesis, initial Einstein--QED record through Round 12.
% Owned by the early-history researcher.  No accepted repository files modified.
\section{From semiclassical backreaction to a controlled finite gauge theory}
\label{early:section}

The early program asked how a quantum current, electromagnetic backreaction,
and gravitational constraints could be calculated consistently.  The surviving
record begins with Round 3.  It contains imported prescribed-background results
and descriptions of earlier research fronts, but no independently identifiable
Round 1 or Round 2 package.  We therefore do not reconstruct their missing
contracts or assign them new results.  The present section accounts for all
31 named studies in Rounds 3--12, using the Round 25 history inventory and the
underlying mathematical reports.  These studies precede the explicit paired
loop contracts introduced later.  Here \workstar\ marks a contribution to this
workbench, including an implementation or an independently derived obstruction;
it does not establish priority over the mathematical literature.

\subsection{The original closure problem and its finite electromagnetic realization}
\label{early:qed}

\paragraph{Initial scope.}
The organizing Einstein--QED action contained Einstein gravity, Maxwell and
Dirac terms, possible scalar-dependent coefficients, and curvature counterterms.
After integrating out the electron, the intended causal object was a
closed-time-path functional
\(\Gamma[g_+,A_+;g_-,A_-;\rho_0]\).  Varying one branch before identifying the
branches defines state-dependent current and stress.  This is a specification
of a problem, not its solution.  A local vacuum-persistence rate is not a causal
current; an Euler--Heisenberg expansion cannot be added independently to the
same full electron determinant.  The four contextual fronts---charged de Sitter
discharge, magnetized pair creation, Cauchy-horizon dynamics, and field mixing---
have different geometries and boundary data.  No theorem identifies their
solutions.  Sources: \repo{evidence/qeg-research/round3/chapters/01_master_action.md}
and \repo{evidence/qeg-research/round3/chapters/02_fronts.md}.

\paragraph{Matched finite Maxwell--Dirac closure \workstar.}
Self-consistent field/mode initial-value formulations and adiabatic
regularization have established precedents \cite{kluger1993pair}.
The executable reduction fixes a flat homogeneous geometry, a static magnetic
field, finitely many Landau levels, and a finite positive momentum quadrature.
This early reduction uses Heaviside--Lorentz natural units, $c=\hbar=1$;
later gauge-graph comparisons restore their explicitly stated physical scales.
With \(s=mt\), \(a=eA_z/m\), \(x=eE/m^2\), and \(b=|eB|/m^2\), set
\(p_i=k_i-a\), \(M_i^2=1+2bn_i\), \(\omega_i^2=M_i^2+p_i^2\), and
\(h_i=(M_i,0,p_i)\).  The positive weights are
\(w_i=b(2-\delta_{n_i0})w_{k_i}/(4\pi^2)\).  A prime in this subsection means
\(\dd/\dd s\).  Bloch evolution and three finite sums are
\begin{align}
 a'&=-x,& r_i'&=2h_i\times r_i,\notag\\
 S&=\sum_iw_i(r_{iz}+p_i/\omega_i),&
 C&=\sum_i\frac{w_iM_i^2}{4\omega_i^5},&
 D&=\sum_i\frac{5w_iM_i^2p_i}{8\omega_i^7}.
 \label{early:bloch}
\end{align}
The fixed magnetic matching coefficient and divided field equation are
\begin{align}
 \chi_b&=\frac{e^2}{12\pi^2}
 \left[b-\log(2b)-\psi\left(1+\frac1{2b}\right)\right],\notag\\
 Z&=1+\chi_b-e^2C,& Zx'&=F-e^2(S+Dx^2).
 \qquad\text{\workstar}\label{early:matched}
\end{align}
The susceptibility comparison uses the constant-field, second-order
electric-probe and all-orders magnetic-field setting of
\cite{mages2010euler}; the evolving current is obtained from the retained modes.
Here \(F\) is the prescribed source.  The mode subtraction is
\(c_i=M_i^2x'/(4\omega_i^5)-5M_i^2p_ix^2/(8\omega_i^7)\), with matching
restored once.  It follows by solving
\(2h_i\times r_i^{(2)}=(r_i^{(1)})'\), with
\(r_i^{(1)}=(0,-M_ix/(2\omega_i^3),0)\), and fixing the longitudinal
component by \(|r_i|=1\); the companion energy subtraction
\(M_i^2x^2/(8\omega_i^5)\) differentiates to \(xc_i\).
The finite matching is the difference between the full-momentum magnetic
Landau sum and the zero-field coefficient
\(e^2\log(1+2bN)/(12\pi^2)\), followed by \(N\to\infty\).
In particular \(S+Dx^2\) is a numerator, whereas the full
renormalized matter current in these conventions is
\(J=S-Cx'+Dx^2+(\chi_b/e^2)x'\).  Confusing the two changes the response.

Define \(U=\sum_iw_i(h_i\cdot r_i+\omega_i)\).  Direct differentiation uses
\(p_i'=x\), \(h_i\cdot(h_i\times r_i)=0\), and
\(C'=-2Dx\) to give \(U'=xS\), \(Z'=2e^2Dx\), and therefore
\begin{equation}
 W=\tfrac12 Zx^2+e^2U,
 \qquad W'=xF.\qquad\text{\workstar}
 \label{early:work}
\end{equation}
This is an exact finite-model identity, not a continuum error estimate.
For physical \(|r_i|\le1\), \(U\ge0\).  If \(Z(a)\ge z_*>0\) for all
\(a\), then \(W\ge z_*x^2/2\) and, by applying the differential inequality
first to \(\sqrt{W+\varepsilon}\) and then letting \(\varepsilon\downarrow0\),
\begin{equation}
 \sqrt{W(s)}\le\sqrt{W(0)}+
 \frac{1}{\sqrt{2z_*}}\int_0^s|F(u)|\,\dd u.
 \label{early:existence}
\end{equation}
The spins remain in compact balls, \(x\) is bounded on each finite interval,
and \(a=a_0-\int x\) cannot diverge there.  Local Lipschitz continuity of the
finite vector field proves global classical continuation for continuous
drives, as assumed in the historical theorem.  An initially mixed
mode satisfying \(|r_j|\le1-\delta\) also gives
\(U\ge w_j\delta|k_j-a|\), hence bounds \(a\) for an integrable drive.
That additional premise does not hold for every pure-vacuum production run.

The historical source-off calculation reached \(x(50)\simeq-0.03757\) with
4096 momentum nodes, Landau levels 0--8, and window \(K=40\).  The recorded
2048-to-4096 history difference \(3.26\times10^{-9}\) and separate Landau
endpoint difference \(5.06\times10^{-8}\) are refinement diagnostics.
The advisor's \(10^{-18}\) absolute benchmark failed at approximately
\(1.82\times10^{-18}\); it remains failed.  Independent spinor/Bloch and
matching controls established a further methodological point: a consistently
wrong prescription may conserve its own energy.  The separate massless
lowest-Landau-level comparator obeys
\(x''+\Omega_B^2x=F'\), \(\Omega_B^2=e^2b/(2\pi^2)\), in its own declared
scaling; it is not obtained by dividing the massive variables by zero.
Sources: \repo{evidence/qeg-research/round3/theory.md},
\repo{evidence/qeg-research/round3/README.md}, and
\repo{evidence/qeg-research/round5/theorem_advisor.md}.

For the weak probe about a stationary vacuum, eliminating the transverse
components of the linearized Bloch equation supplies the oscillator denominator
\(4\omega_i^2-(\nu+i0)^2\).  The same zero-frequency subtraction gives
\begin{equation}
 \chi(b,\nu)=\chi_b+e^2\sum_i\frac{w_iM_i^2}{\omega_i^3}
 \left[\frac1{4\omega_i^2-(\nu+i0)^2}-\frac1{4\omega_i^2}\right].
 \qquad\text{\workstar}\label{early:weakresponse}
\end{equation}
This tests matching independently of the nonlinear work identity.
In an instantaneous eigenbasis, the mode current also contains a coherence:
\(r_{iz}+p_i/\omega_i=2p_if_i/\omega_i-M_iv_i^{\rm coh}/\omega_i\).
Retaining the occupation \(f_i\) alone cannot reproduce that current.

\paragraph{Causal tangent and the quantity a stability theorem controls \workstar.}
Let \(v=\delta a\), \(u=\delta x\), \(q_i=-v\), \(\eta_i=\delta r_i\),
and \(f=\delta F\), with grid and matching held fixed.  Differentiation gives
\begin{align}
 \delta S&=\sum_iw_i(\eta_{iz}+M_i^2q_i/\omega_i^3),\notag\\
 \delta C&=-\sum_i\frac{5w_iM_i^2p_iq_i}{4\omega_i^7},&
 \delta D&=\sum_i\frac{5w_iM_i^2(M_i^2-6p_i^2)q_i}{8\omega_i^9},\notag\\
 v'&=-u,&
 u'&=\frac{f-e^2(\delta S+\delta D x^2+2Dxu)+e^2\delta Cx'}{Z},\notag\\
 \eta_i'&=2h_i\times\eta_i+2q_i e_z\times r_i.
 \qquad\text{\workstar}\label{early:tangent}
\end{align}
The fundamental-matrix integral gives causal support and finite-time
Gronwall bounds.  The first variation \(\delta W\), satisfying
\((\delta W)'=uF+xf\), is signed and cannot be used as a positive norm.
At a stationary pure vacuum, restrict to variations preserving the Bloch
norm, so \(h_i\cdot\eta_i=0\), and define
\(t_i=\partial_{p_i}(h_i/\omega_i)\) and
\(\xi_i=\eta_i+q_it_i\).  Then
\(Z_0u'=f-e^2\sum_iw_i\xi_{iz}\),
\(\xi_i'=2h_i\times\xi_i+ut_i\), and
\begin{equation}
 E_2=Z_0u^2+e^2\sum_iw_i\omega_i|\xi_i|^2,
 \qquad E_2'=2uf.\qquad\text{\workstar}
 \label{early:secondvariation}
\end{equation}
Indeed \(h_i\cdot\xi_i=0\) and \(\omega_i\xi_i\cdot t_i=\xi_{iz}\)
cancel the exchange terms.  This controls \(u,\xi\), not every raw tangent.
For one mode with \(M=w=e^2=1\), \(p_0=0\), and \(Z_0=3/4\),
\(u=1\), \(v=-s\), \(\xi=(0,-1/2,0)\),
\(\eta=(0,-1/2,-s)\) solve the unforced tangent while \(E_2=1\).
The secular raw tangent defeats a stronger stability claim.  Rounds 4 and 6
also retained source-center and acceptance repairs rather than overwriting
earlier failures.  Sources: \repo{evidence/qeg-research/round4/response_contract.md}
and \repo{evidence/qeg-research/round6/extensions.md}.

\paragraph{Finite-window positivity and its cofinal obstruction \workstar.}
For continuous momentum in \([-K,K]\), write
\begin{align}
 C_{N,K}(a)&=\frac{b}{16\pi^2}\sum_{n=0}^Nd_n
 [P_{M_n}(K-a)+P_{M_n}(K+a)],\notag\\
 P_M(y)&=\frac{u-u^3/3}{M^2},\qquad u=\frac{y}{\sqrt{M^2+y^2}}.
 \qquad\text{\workstar}\label{early:window}
\end{align}
Since \(P_M'=M^2/(M^2+y^2)^{5/2}\) is even and decreases with \(|y|\),
\(C\) is even and strictly decreases for \(a>0\).  Its maximum is at zero.
This statement is about an integral: a two-node Gauss rule with nodes
\(\pm K/\sqrt3\) has, at a node, a lowest-level contribution
\(e^2C_{\rm disc}\ge\alpha_{\rm em}bK/(4\pi)\), which can exceed the
continuous bound \(\alpha_{\rm em}b/(3\pi)\).
Independently, at \(b=10,K=20\), positivity of every quadrature weight and
exactness on constants, together with
\(\alpha_{\rm em}=e^2/(4\pi)<1/137\), suffice for the rational estimate
\begin{equation}
 e^2C_{\rm disc}<\frac{67743204500}{274510827927}<\frac14,
 \qquad Z>\frac34,
 \qquad\text{\workstar}\label{early:rationalz}
\end{equation}
uniformly in every finite Landau cutoff.  The estimate uses mass bounds for
the first levels and the integral majorant for the decreasing tail.
It verifies a premise of \eqref{early:existence}, without certifying a trajectory.

At infinite momentum window the integral is elementary:
\begin{equation}
 C_{N,\infty}=\frac{b}{12\pi^2}
 \left(1+2\sum_{n=1}^N\frac1{1+2bn}\right),\qquad
 Z_{N,\infty}=1-\frac{e^2}{12\pi^2}
 \left[\log(2b)+\psi\left(N+1+\frac1{2b}\right)\right].
 \qquad\text{\workstar}\label{early:cofinal}
\end{equation}
The latter tends to \(-\infty\).  Positivity makes \(C_{N,K}(0)\) increase
in both cutoffs: choose a finite \(N\) exceeding any prescribed threshold,
then a finite \(K\) approximating that integral.  Every cofinal cutoff path
eventually exceeds both.  Thus unchanged matching cannot retain a positive
uniform coefficient margin.  This obstructs this divided-coefficient proof
route, not continuum QED or every renormalization prescription.  A per-run
reference choice \(\chi_{N,K}=z_{\rm ref}-1+e^2C_{N,K}(0)\) yields
\(Z=z_{\rm ref}+e^2[C(0)-C(a)]\ge z_{\rm ref}\) for the continuous window;
it changes the matching and does not inherit that inequality for arbitrary
quadrature.  A scalar contribution uniformly bounded over the permitted
field domain cannot cancel the cofinal divergence at fixed \(a\), or uniformly
bounded \(a\).  Sources: \repo{evidence/qeg-research/round6/extensions.md} and
\repo{evidence/qeg-research/round7/renormalization_map.md}.

The same estimates give the coefficient-only tail
\(e^2\sup_a|C_{\infty,K}-C_{N,K}|\le
\alpha_{\rm em}K/[\pi\sqrt{1+2bN}]\), provided previously retained
quadratures are unchanged.  This does not bound omitted dynamical modes.
Whenever a finite cutoff has \(Z(0)<0\), continuity and
\(Z(a)\to1+\chi_b>0\) produce singular coefficient surfaces at opposite
values of \(a\); no argument here proves that a physical trajectory reaches
one.

\subsection{What gravitational and scalar extensions actually add}
\label{early:gravity}

\paragraph{An anisotropic current--stress interface \workstar.}
For \(ds^2=-dt^2+a_\perp^2(dx^2+dy^2)+a_\parallel^2dz^2\), define
\(h=\dot a_\perp/a_\perp\), \(k=\dot a_\parallel/a_\parallel\), and
\(\theta=2h+k\).  Maxwell evolution gives
\(\dot E+2hE=-j\), \(\dot B+2hB=0\), with field stresses
\((\rho,p_\perp,p_\parallel)_{\rm EM}=(u,u,-u)\),
\(u=(E^2+B^2)/2\).  The spatial Einstein equations
\(\dot h=(\Lambda-\kappa p_\parallel-3h^2)/2\) and
\(\dot k=\Lambda-\kappa p_\perp-h^2-k^2-hk-\dot h\) imply, by differentiation,
\begin{align}
 C_g&=h^2+2hk-\kappa\rho-\Lambda,&
 Q&=\dot\rho+2h(\rho+p_\perp)+k(\rho+p_\parallel),\notag\\
 \dot C_g&=-\theta C_g-\kappa Q,&
 C_g(t)&=\frac{V(0)C_g(0)-\kappa\int_0^tV(s)Q(s)\,\dd s}{V(t)},
 \quad V=a_\perp^2a_\parallel.
 \qquad\text{\workstar}\label{early:constraint}
\end{align}
The coefficient is \(-\theta\) for these off-constraint equations, not the
\(-2\theta\) initially suggested in review.  Running \(\kappa,\Lambda\)
adds \(-\dot\kappa\rho-\dot\Lambda\).  Moreover
\(\dot(k-h)+\theta(k-h)=\kappa(p_\parallel-p_\perp)\).
The replacement \((\delta p_\perp,\delta p_\parallel)=(kf,-2hf)\) leaves
the energy equation unchanged but changes the shear source by
\(-\kappa\theta f\).  Exact mean-energy conservation therefore cannot identify
both pressures.

At a fixed comoving regulator the rescaled Dirac modes have
\(H_n=m\beta+P\alpha_z+c_n\alpha_x\),
\(P=(k_z-qA_z)/a_\parallel\), \(c_n=\sqrt{2n|qB_0|}/a_\perp\).
Writing \(\mathcal D=|qB_0|/(4\pi^2V)\) and \(\mathcal S[X]\) for the
occupied-mode trace and sum, variations of this same Hamiltonian give
\begin{equation}
 (\rho_m,j,p_\parallel,p_\perp)=\mathcal D
 \big(\mathcal S[H],q\mathcal S[\alpha_z],
 \mathcal S[P\alpha_z],\tfrac12\mathcal S[c_n\alpha_x]\big).
 \qquad\text{\workstar}\label{early:pressures}
\end{equation}
Using \(\dot{\mathcal D}=-\theta\mathcal D\),
\(\dot P=qE-kP\), and \(\dot c_n=-hc_n\) proves
\(\dot\rho_m+2h(\rho_m+p_\perp)+k(\rho_m+p_\parallel)=Ej\).
The two occupied columns at \(n\ge1\) already include the spin multiplicity.
A rotation of \(P\sigma_3+m\sigma_1-sc_n\sigma_2\) to real transverse mass
adds the connection
\(s(m\dot c_n-c_n\dot m)/[2(m^2+c_n^2)]\) to \(P\); omitting it explains
why the naive effective-mass replacement fails in the imported expanding
background comparator.  These are bare regulated quantities, not renormalized
Einstein sources.

A finite counterterm \(-c\int\sqrt{-g}F^2/4\) changes the current by
\(c(\dot E+2hE)\) and the stress by \(c(u,u,-u)\); its Ward exchange cancels
only when both are changed.  The classical Ohmic and later scalar-gravity
fixtures test this interface with their own matter.  They do not compute a
common covariant quantum current and directional stress.  The archived physical
electron scale estimate \(\kappa\rho_{\rm EM}/m_e^2\simeq2.40\times10^{-39}\)
at \(E/E_c=1,B/B_c=100\) also rules out interpreting the fixture's visibly
large gravitational response as an electron-coupling prediction.  Sources:
\repo{evidence/qeg-research/round3/gravity.md},
\repo{evidence/qeg-research/round3/evidence/prior_curved_sector_results.json},
and \repo{evidence/qeg-research/round5/research_bridge.md}.
The locally covariant Dirac analysis \cite{zahn2014dirac} provides a
renormalization framework with explicit background qualifications; its direct
free stress-conservation discussion does not supply the strong-field
Bianchi-I mode sums required here.

Two retained logical controls distinguish conservation from validity.
A random energy equal to \(\bar E/p\) with probability \(p\), and zero
otherwise, has the same mean as the eigenvalue \(\bar E\) but variance
\(\bar E^2(p^{-1}-1)\).  Mean stress alone therefore does not determine
stress noise.  Also exact local evolution
\(e^{-i\eta\sigma_x}\otimes e^{-i\eta\sigma_x}|00\rangle\) stays separable,
whereas normalizing its first-order truncation gives
\begin{equation}
 \operatorname{Var}(E)=\bar E^2(p^{-1}-1),\qquad
 \mathcal C_{\rm truncated}=\frac{2\eta^2}{1+2\eta^2},\qquad
 \mathcal C_{\rm exact}=0.\qquad\text{\workstar}
 \label{early:logicalcontrols}
\end{equation}
The latter apparent witness is at the discarded order.
These project counterexamples \workstar\ do not diagnose gravitational
breakdown or a physical entanglement experiment.  For the inherited
Cauchy-horizon profile \(\phi\sim|V|^\beta\), integrating
\(|\partial_V\phi|^2\sim|V|^{2\beta-2}\) gives the familiar threshold
\(\beta>1/2\), with logarithmic divergence at equality; it is an energy
regularity test, not a solved backreacted horizon.

\paragraph{Action-derived scalar completion \workstar.}
A dynamical coefficient requires its own equation and energy.  Round 7 chose
\(f(\phi)=1+g^2\phi^2\), scalar kinetic energy \(y^2/2\), and
\(V(\phi)=\nu^2\phi^2/2\), with \(y=\phi'\).  Variation of
\(-f(\phi)F^2/4\), with the finite-mode current unchanged, gives
\begin{align}
 Zx'&=F-e^2(S+Dx^2)-f_\phi yx,&
 y'&=-\nu^2\phi+\tfrac12f_\phi(x^2-b^2),\notag\\
 \widetilde W&=\tfrac12Zx^2+e^2U+\tfrac12y^2+
 \tfrac12(\nu^2+g^2b^2)\phi^2,&
 \widetilde W'&=xF.
 \qquad\text{\workstar}\label{early:scalar}
\end{align}
The scalar and field exchange terms cancel directly.  Since \(f\ge1\), the
finite \(Z>3/4\) certificate is retained, and the same continuation argument
applies; the uncoupled zero-frequency scalar is handled as an affine solution.
By contrast, merely prescribing \(\chi(s)\) in the old equations produces the
unaccounted power \(\chi'x^2/2\).  The separate classical Bianchi-I model
with \(f=e^{2c\varphi}\) closes its own scalar/Maxwell exchanges.  It supplies
no identification with quantum stress.  Renormalization-scale running,
cutoff changes, and a dynamical scalar are different operations.  Sources:
\repo{evidence/qeg-research/round7/research_report.md} and
\repo{evidence/qeg-research/round7/renormalization_map.md}.

\subsection{Spectral logic before a nonabelian calculation}
\label{early:spectral}

\paragraph{False plateaux and conditional spectral bridges \workstar.}
The move to pure Yang--Mills changes the theory and its target
\cite{jaffe_witten}.  Its useful
inheritance is explicit control of assumptions, not an abelian mass-gap result.
For a nonzero finite positive spectral measure supported on \([0,\infty)\),
\(C(t)=\int e^{-Et}\,\dd\mu(E)\), let
\(E_*=\inf\operatorname{supp}\mu\).  Then
\(C(t+h)\le e^{-hE_*}C(t)\); logarithmic convexity follows from H\"older.
Consequently the effective mass
\(m_h(t)=h^{-1}\log[C(t)/C(t+h)]\) decreases to \(E_*\) from above.
For \(t>0\), differentiating gives
\begin{equation}
 -\partial_t\log C=\langle E\rangle_t,
 \qquad \partial_t\langle E\rangle_t=-\operatorname{Var}_t(E).
 \label{early:effective}
\end{equation}
Positive measure in every \([E_*,E_*+\varepsilon]\) fixes the asymptotic
exponential rate; convexity then fixes the secant limit.  Thus an early plateau
near 1 in the executed \(e^{-t}+10^{-12}e^{-0.1t}\) example is not a lower
bound on the threshold 0.1.  A lighter orthogonal channel may be invisible
even at infinite time.  The periodic massless free-field diagnostic explicitly
excluded its zero mode; its lowest nonzero momentum is not the full gap.

Two standard spectral-theorem arguments were made explicit.  First, if
\(H_n\to H\) strongly in resolvent on a common Hilbert space and all
\(\operatorname{spec}(H_n)\cap(0,m)=\varnothing\) for the same \(m>0\),
then compactly supported continuous \(f\) in that interval satisfies
\(f(H)=\lim f(H_n)=0\).  Positive but nonuniform gaps are insufficient:
\(H_ne_k=\max(1/n,1/k)e_k\), \(k\ge1\), and \(H_ne_0=0\), converge in
norm to a gapless diagonal operator.  Second, if a reconstructed positive
Hamiltonian already exists and centered states \(\psi_O\) span densely in
\(\Omega^\perp\), a common bound
\(\langle\psi_O,e^{-tH}\psi_O\rangle\le A_Oe^{-mt}\) at all sufficiently
large times implies
\begin{equation}
 e^{-(m-\varepsilon)t}\norm{P_H([0,m-\varepsilon])\psi_O}^2
 \le A_Oe^{-mt}\quad\Longrightarrow\quad
 H|_{\Omega^\perp}\ge m.
 \label{early:densebridge}
\end{equation}
Density removes every lower channel.  These known implications do not construct
the continuum Hilbert space or supply the common physical rate.  The Round 5
and Round 8 planners replay finite reviewed implication libraries; they are
not proof-assistant kernels.  The Round 8 delayed-source refinement met its
declared adjacent-grid working scale, while 21 symbolic checks were blocked
by an unavailable dependency and remain unexecuted in that audit.  Sources:
\repo{research/round8/spectral-proof.md},
\repo{research/round8/advisor-dossier.md}, and \repo{research/round8/README.md}.

\paragraph{Finite Haar hierarchy, and what a single identity cannot certify \workstar.}
For the finite Wilson measure \(\dd\mu_W=Z^{-1}e^{-S_W}\dd H\) with
\(S_W=\beta_W\sum_p(1-\Tr U_p/2)\), left Haar differentiation yields
\(\langle L_af-fL_aS_W\rangle=0\).  Conversely, on a finite product of
connected compact groups, if every smooth insertion satisfies every such
identity under a probability measure \(\nu\), substitute \(f=e^S\varphi\).
Then \(\int L_a\varphi\,e^S\dd\nu=0\).  Integrating the generated flows
shows \(e^S\nu\) is translation invariant; uniqueness of Haar measure gives
\(\nu=Z^{-1}e^{-S}H\), without assuming absolute continuity in advance.

For an independent link and an unnormalized staple \(A\), write
\(UA=a_0I+i\boldsymbol a\cdot\boldsymbol\sigma\).  The generator
\(\exp(it\sigma_a/2)\) gives \(L_aS_W=\beta_Wa_a/2\) and
\(L_a^2S_W=\beta_Wa_0/4\), hence
\begin{equation}
 \left\langle\beta_W^2|\boldsymbol a|^2-3\beta_Wa_0\right\rangle=0.
 \label{early:ward}
\end{equation}
The staple is not a unit quaternion.  The identity assumes the varied link
occurs once in each incident plaquette.  On one SU(2) variable the wrong law
\((\delta_I+\delta_{-I})/2\) passes this residual exactly but gives
\(\langle L_a(u_0u_a)-u_0u_aL_aS\rangle=1/2\).
The complete hierarchy characterizes the law; one measured residual does not.
The six finite four-dimensional lattice chains at \(\beta_W=0,0.5,2.2\)
tested the declared nonabelian sampler and local identities.  The
\(\beta_W=2.2\) autocorrelation evidence was insufficient.  No chain supplied
a continuum gap.  Sources: \repo{research/round9/advisor/advisor.md} and
\repo{research/round9/lattice/RESULTS.md}.

\paragraph{An exact auxiliary transfer operator and the missing norm \workstar.}
For central convolution by \(e^{\beta\cos\theta}/Z_\beta\) on full
\(L^2(SU(2))\), the class Haar measure is
\((2/\pi)\sin^2\theta\,\dd\theta\).  With \(n=2j+1\), character integration
gives \(c_j=I_{n-1}-I_{n+1}=2nI_n/\beta\); Schur orthogonality divides by
\(n\), yielding the complete spectrum
\begin{equation}
 r_j=\frac{I_{2j+1}(\beta)}{I_1(\beta)},\quad
 \operatorname{mult}(r_j)=(2j+1)^2,\quad
 \Delta=-a_t^{-1}\log\frac{I_2(\beta)}{I_1(\beta)}.
 \label{early:convolution}
\end{equation}
For \(\beta>0\), the positive integral representation identifies
\(I_{n+1}/I_n\) with a mean of \(t\in(-1,1)\) tilted by \(e^{\beta t}\);
pairing \(t,-t\) proves this mean is strictly between zero and one.  The
constant is the unique stationary state.  At fixed spin,
\(-\log r_j=2j(j+1)/\beta+O(\beta^{-2})\).  If
\(\beta a_t\to c>0\), dominated convergence of square-summable Peter--Weyl
coefficients gives \(T_\beta^{\lfloor t/a_t\rfloor}\to e^{-2t\mathcal C/c}\)
strongly, with rotor gap \(3/(2c)\).  Other choices of the clock produce zero
or infinite limiting gaps.  Independent endpoint gauge invariance of an open
link forces constants, so this rotor is not the interacting gauge transfer
operator.  Negative \(\beta\) supplies a pointwise positive kernel with a
negative fundamental eigenvalue, refuting positivity by inspection alone.

An additional exact counterexample isolates an operator-norm error.  On
\(N\) uniformly weighted atoms let \(B\) average the first two coordinates
and fix the rest, \(\Pi\) average all coordinates, and
\(v=(1,-1,0,\ldots)/\sqrt2\).  Set
\(P_0=(1-c)B+c\Pi\), \(P_1=P_0+\delta vv^T\), with
\(0<\delta<1-c<1\).  Both are symmetric positive Markov matrices with
strictly positive entries, yet for their density kernels \(K_i=NP_i\),
\begin{equation}
 \norm{P_1-P_0}_{2\to2}=\delta,
 \qquad \norm{K_1-K_0}_{L^1(H\otimes H)}=\frac{2\delta}{N}.
 \qquad\text{\workstar}\label{early:normcounter}
\end{equation}
Thus product-space \(L^1\) control cannot replace uniform row and column
control in Schur's test.  For actual positive contractions sharing a vacuum,
\(\norm{T|_{\Omega^\perp}}\le e^{-ma}\) and
\(\norm{S-T}\le\varepsilon_a\) imply the conditional gap
\(-a^{-1}\log(e^{-ma}+\varepsilon_a)\), when its argument is below one.
Telescoping gives \(\norm{S^n-T^n}\le n\varepsilon_a\); fixed-time convergence
needs \(\varepsilon_a=o(a)\) by this route.  Sources:
\repo{research/round9/transfer/PROOF.md} and
\repo{research/round9/advisor/advisor.md}.

\subsection{Gauge reduction, infinite representation tails, and certified dynamics}
\label{early:gauge}

\paragraph{One-square certification \workstar.}
For an open four-link square, all vertex Gauss laws reduce the physical space
to class functions on SU(2).  Assume \(\alpha>0\), \(\lambda\ge0\).
The known character and radial reductions give
\begin{equation}
 H_{nm}=[\alpha n(n+2)+\lambda]\delta_{nm}
 -\tfrac\lambda2(\delta_{n,m+1}+\delta_{n,m-1}),\quad n,m\ge0,
 \label{early:jacobi}
\end{equation}
The four link Casimirs contribute \(4j(j+1)=n(n+2)\), while
\(2x\chi_{n/2}=\chi_{(n+1)/2}+\chi_{(n-1)/2}\), with the negative-index
character absent, fixes the off-diagonal factor.
Equivalently, the Hamiltonian is
\(-\alpha\partial_\theta^2-\alpha+\lambda(1-\cos\theta)\)
on Dirichlet functions \(u=\sqrt{2/\pi}\sin\theta f\) on \((0,\pi)\).
The Mathieu representation is established prior work \cite{bauer2023basis}; the workbench's
contribution is the explicit certified error contract.  Compact resolvent
and uniqueness for the second-order Dirichlet equation imply a simple ground
state and a positive gap for fixed \(\alpha>0\).  Retain \(n<N\), let
\(A=P_NHP_N\), \(\tau=\alpha N(N+2)\), and choose a certified
\(\mu_1(A)<R<\tau\).  The sole cross-boundary entry is \(-\lambda/2\).
Young's inequality and tail positivity give
\begin{equation}
 B=A-\frac{(\lambda/2)^2}{\tau-R}|N-1\rangle\langle N-1|,
 \qquad \mu_k(B)\le E_k(H)\le\mu_k(A),\quad k=0,1.
 \qquad\text{\workstar}\label{early:onetail}
\end{equation}
Exact rational Sturm inertia, including zero and endpoint policies, encloses
these eigenvalues.  After dropping the scalar \(\lambda I\), bounded
multiplication proves \(|\delta(\kappa)-\delta(\kappa')|\le2|\kappa-\kappa'|\)
for \(\delta=\Delta/\alpha\), \(\kappa=\lambda/\alpha\).  A certified interval
cover then gave \(\Delta\ge0.999999\alpha\) on \(0\le\kappa\le10\).
Point samples alone would not prove this.  A two-dimensional trial space
localized in \(0<\theta<w\) gives
\(\Delta\le4\pi^2\alpha/w^2+\lambda w^2/2-\alpha\), hence
\(\Delta\le2\sqrt2\pi\sqrt{\alpha\lambda}-\alpha\to0\) at fixed
\(\lambda>0\) as \(\alpha\downarrow0\), with the chosen
\(w=(8\pi^2\alpha/\lambda)^{1/4}\le\pi\) once
\(\alpha/\lambda\le\pi^2/8\).  This path is not a matched
continuum trajectory.  Source: \repo{research/round10/advisor/advisor.md}.

\paragraph{The interacting two-square quotient and its complete electric spectrum \workstar.}
The canonical Hamiltonian framework is established \cite{kogut1975hamiltonian}.
The seven-link, six-vertex open graph has physical space
\(L^2(SU(2)^2)^{\mathrm{Ad,diag}}\).  Tree gauge leaves loops \(U,V\) based
at the shared vertex.  This planar spatial patch with continuous time is a
finite \(2+1\)-dimensional regulator; increasing its polynomial cutoff neither
adds spatial cells nor produces the four-dimensional target.
The complete orbit coordinates are
\(x=\Tr U/2\), \(y=\Tr V/2\), \(z=\Tr(UV)/2\), with
\begin{equation}
 \Omega=\{ |x|,|y|\le1:\;1-x^2-y^2-z^2+2xyz\ge0\},
 \qquad\dd\mu=\frac2{\pi^2}\dd x\,\dd y\,\dd z.
 \qquad\text{\workstar}\label{early:quotient}
\end{equation}
To derive it, write the SU(2) quaternions as \((x,\boldsymbol a)\) and
\((y,\boldsymbol b)\).  The three traces fix their Gram matrix, and two
vector pairs with that Gram matrix differ by a proper rotation.  The scalar
marginals are semicircular and the relative angle has uniform cosine;
the Jacobian to \(z=xy-\boldsymbol a\cdot\boldsymbol b\) cancels both square
roots.  This proves the measure as well as orbit classification, including
central pairs.  The pair \((x,y)\) alone loses the relative orientation.

The shared derivative acts as \(L_U-R_V\), so the kinetic operator is
\(K=3(C_U+C_V)-\sum_a(L_U^a-R_V^a)^2\), rather than two independent rotors.
In trace coordinates,
\begin{equation}
 K=-\operatorname{div}(A\nabla),\quad
 A=\begin{pmatrix}
 1-x^2&(z-xy)/4&3(y-xz)/4\\
 (z-xy)/4&1-y^2&3(x-yz)/4\\
 3(y-xz)/4&3(x-yz)/4&3(1-z^2)/2
 \end{pmatrix}.
 \qquad\text{\workstar}\label{early:operator}
\end{equation}
Its domain is inherited from invariant \(H^2(SU(2)^2)\), equivalently the
closure of the polynomial core, not Dirichlet data on the quotient boundary.
The checks \(Kz=9z/2\) and \(K(xy)=13xy/2-z/2\) detect the shared link.
For a monomial \(x^ay^bz^c\), its degree-preserving coefficient is
\begin{equation}
 \epsilon_{abc}=a^2+b^2+\tfrac32c^2+\tfrac12ab+
 \tfrac32c(a+b)+2a+2b+3c,
 \qquad m_d=\tfrac58d^2+2d+\tfrac38(d\bmod2).
 \qquad\text{\workstar}\label{early:shell}
\end{equation}
All other terms lower degree.  Symmetry makes each orthogonal polynomial
shell reduce \(K\); the diagonal quotient coefficients are its eigenvalues.
Density of polynomials proves completeness.  Writing
\(j=(a+c)/2,k=(b+c)/2,\ell=(a+b)/2\) gives
\(\epsilon=3j(j+1)+3k(k+1)+\ell(\ell+1)\), independently identifying the
spin-network spectrum.  Its low electric levels agree with the independent
two-plaquette construction in \cite{froland2025two}; this is a comparison of
those levels, not an identification of all intermediate gauge-fixed formulas.
At fixed \(d=a+b+c\), the minimum has \(c=0\) and
\(a,b\) as equal as possible, giving \(m_d\).  Thus a degree cutoff \(D\)
has dimension \(\binom{D+3}{3}\) and exact omitted threshold \(m_{D+1}\).
Sources: \repo{research/round10/skeptic/coordinate-note.md} and
\repo{research/round11/advisor/advisor.md}.

\paragraph{Interacting gap enclosures and the cost of a global bound \workstar.}
For \(H=\alpha K+\lambda_1(1-x)+\lambda_2(1-y)\), with
\(\alpha>0\), \(\lambda_1,\lambda_2\ge0\), bounded perturbation and
positivity improvement on the connected compact group give compact resolvent
and a unique positive physical ground state.  With \(P=P_D\),
\(M=-\lambda_1x-\lambda_2y\), and \(\tau=\alpha m_{D+1}\), the exact
boundary form is finite:
\begin{equation}
 \begin{aligned}
 CC^*&=PM(P_{D+1}-P_D)MP,\qquad
 B=PHP-\frac{CC^*}{\tau-R},\\
 \mu_k(B)&\le E_k(H)\le\mu_k(PHP),\qquad k=0,1.
 \qquad\text{\workstar}
 \end{aligned}\label{early:twotail}
\end{equation}
Here \(\mu_1(PHP)<R<\tau\).  Multiplication raises degree by at most one;
Young's inequality then proves the same form comparison as
\eqref{early:onetail}.  In raw monomials, inertia must be computed for
\(H_{\rm form}-rG\), using the Haar Gram matrix \(G\).  Its entries are exact:
if \(\mu_{2r}=\mathrm{Catalan}(r)/4^r\), \(\mu_{2r+1}=0\), and
\(J(n,h)=\sum_{q=0}^h(-1)^q\binom hq\mu_{n+2q}\), then
\begin{equation}
 \langle x^ay^bz^c\rangle=
 \sum_{h=0}^{\lfloor c/2\rfloor}\frac{\binom c{2h}}{2h+1}
 J(a+c-2h,h)J(b+c-2h,h).
 \label{early:moments}
\end{equation}
This follows from the same relative-angle integration.  The gap is
2-Lipschitz in the \(\ell^1\) distance of \((\lambda_1/\alpha,\lambda_2/\alpha)\).
Exact cover certificates gave a floor exceeding \(0.96\alpha\) on
\([0,2]^2\), with recorded lower endpoint \(0.963936567779819\alpha\).
At \((\alpha,\lambda_1,\lambda_2,\rho)=(1,2,3,1)\), the recorded full-space
gap interval was \([3.362386061099721,3.362408390681595]\).

An analytic estimate exposed the volume issue before larger computation.
The SU(2) heat kernel at Casimir time 4 lies between \(3/4\) and \(5/4\):
its nonconstant character sum is bounded by
\(4e^{-3}+\sum_{n\ge2}(2/25)^n<1/4\).
For the two-square kinetic operator, time \(4/3\) therefore gives kernel
bounds \(9/16,25/16\).  Feynman--Kac bounds a positive ground state's ratio
by \((25/9)e^{8(\lambda_1+\lambda_2)/(3\alpha)}\).  The ground-state transform
and free physical Poincar\'e constant 3 give
\begin{equation}
 \Delta\ge\frac{243}{625}\alpha
 e^{-16(\lambda_1+\lambda_2)/(3\alpha)}.
 \qquad\text{\workstar}\label{early:globalgap}
\end{equation}
On any finite connected open hypercubic graph containing a square, the free
physical gap is exactly \(3\alpha\): a nonzero spin support has no degree-one
vertex, hence contains at least four edges, each costing at least \(3/4\);
a fundamental square attains the bound.  Repeating the global heat comparison
for \(E\) links gives
\(\Delta_\Gamma\ge3\alpha(3/5)^{2E}e^{-16\sum_p\lambda_p/\alpha}\).
Its deterioration is not a proof of gap closure.  For \(q\) disconnected
copies, tensor eigenvectors give the exact same first gap as one copy while
the global density-ratio bound decays exponentially in \(q\).

There is also a direct measure obstruction.  For the one-square Hamiltonian,
\begin{equation}
 \frac{[4C+\kappa(1-x)]e^{\theta x}}{e^{\theta x}}
 =\kappa-\theta^2+(3\theta-\kappa)x+\theta^2x^2.
 \qquad\text{\workstar}\label{early:gibbs}
\end{equation}
It is constant only for \(\theta=\kappa=0\).  The square root of a nontrivial
Wilson Gibbs density is thus not this Hamiltonian's ground state.  A
Langevin Poincar\'e theorem for the former cannot silently bound the latter;
the fixed-spacing strong-coupling results of \cite{shen2022strong} retain their
own measure, normalization, and generator hypotheses.
For a declared Hamiltonian convention,
\(\alpha=g_H^2/(2a)\), \(\lambda=2/(g_H^2a)\), so
\(\lambda/\alpha=4/g_H^4\); neither a fixed graph nor a compact coupling
square follows a growing-volume weak-coupling continuum limit.  Sources:
\repo{research/round11/advisor/advisor.md} and
\repo{research/round12/volume/volume-bridge.md}.

A possible replacement must control the actual ground state.  The conditional
Bakry--\'Emery bound
\(\operatorname{Ric}-2\operatorname{Hess}\log\phi_\Gamma\ge c_*g\), uniformly
in volume, would give \(\Delta_\Gamma\ge\alpha c_*\) through the
ground-state generator \(\Delta+2\nabla\log\phi_\Gamma\cdot\nabla\).
No such bound was established.  Similarly, transferring
\(A_j\ge m_*(I-P_j)\) to a continuum operator needs a common realization,
strong convergence of rank-one vacuum projections, and a limiting semigroup
strongly continuous at zero.  The example \(A_j=j(I-P)\) has
\(e^{-tA_j}\to P\) for each \(t>0\), but no strongly continuous limit on
the full space.  Lower-gap control alone can leave only a vacuum; nontriviality
is a separate obligation.

\paragraph{A full-Hilbert-space dynamical error bound \workstar.}
The earlier driven one- and two-square histories were finite-Galerkin studies.
Round 12 added the omitted representation error.  Let
\(H(t)=\alpha(t)K_\rho+\lambda_1(t)(1-X)+\lambda_2(t)(1-Y)\),
\(K_\rho=3(C_U+C_V)-\rho\sum(L_U-R_V)^2\), with fixed \(\rho>0\).
Remove the common scalar phase.  The electric propagator preserves every
degree shell and the remaining bounded interaction has norm at most
\(w=|\lambda_1|+|\lambda_2|\), with shell bandwidth one.  Real locally
integrable coefficients suffice for mild unitary evolution.
For normalized initial support in \(P_{d_0}\), let \(\psi_D\) be exact
Galerkin evolution and \(A(t)=\int_0^tw(s)\,\dd s\).  Duhamel in the finite
tail block gives inductively
\(\norm{(P_D-P_{n-1})\psi_D}\le A^{n-d_0}/(n-d_0)!\).
The only full-equation defect is
\(r_D=\Pi_{D+1}W\Pi_D\psi_D\).  One more integral yields
\begin{equation}
 \norm{\psi(t)-\psi_D(t)}\le
 \min\left\{2,\frac{A(t)^{D-d_0+1}}{(D-d_0+1)!}\right\}.
 \qquad\text{\workstar}\label{early:factorial}
\end{equation}
No exponential factor is needed because both block propagators are unitary.
Equal action bounds do not identify different protocols.  For a computed
path, add its preparation error and the integral of
\(\norm{i\dot{\widetilde\psi}_D-H_D\widetilde\psi_D}\).
Norm conservation alone does not bound that residual.  The boundary norm
has the exact Gram expression
\(\norm{r_D}^2=c^*[N-MG^{-1}M]c\), where \(M,N\) are the form matrices of
\(P_DWP_D,P_DW^2P_D\).

For a finite step of length \(h\), midpoint \(B=H_D(m)\),
\(C=H_D'(m)\), and \(M_2=\sup\norm{H_D''}\), a Taylor comparison followed
by two unitary Duhamel integrals gives
\begin{equation}
 \norm{U_D-e^{-ihB}}\le h^3M_2/24+h^3\norm{[B,C]}/12
 +h^4\norm C^2/32.
 \qquad\text{\workstar}\label{early:midpoint}
\end{equation}
The first term integrates the quadratic Taylor remainder.  The centered
linear integral cancels its constant interaction-picture part, leaving the
commutator term; the second Duhamel remainder is half the squared integrated
linear norm.  Compressed \(X,Y\) need not commute.  Evaluating the exponential
is a separate error: for a finite selfadjoint \(B\), its order-\(q\) Taylor
remainder is at most \((h\norm B)^{q+1}/(q+1)!\).
The archived exact rational two-step state used degree 3, Taylor order 100,
\((\lambda_1,\lambda_2)=(1/20,1/10)\) then \((1/10,1/20)\), each for
duration 1, and \(\alpha=\rho=1\).  With
\(e=(114/5)^{101}/101!\), its total bound was
\(27/80000+2e+e^2\), including the product error and the unnormalized
computed state.  This certifies that declared protocol, not all historical
drives.  Bounded-observable errors are at most \(2\norm O\) times a normalized
state error; energy needs its own work identity and cannot use a fictitious
global norm of \(H\).  Sources: \repo{research/round12/advisor/advisor.md} and
\repo{research/round12/solver/output/exact_step_certificate.json}.

\paragraph{Closure tests that retain the missing fluctuation \workstar.}
For the one-square Euclidean measure proportional to
\(e^{\kappa x}\sqrt{1-x^2}\,\dd x\), integration of
\(\partial_x[(1-x^2)^{3/2}e^{\kappa x}f]\) gives, with \(m_n=\langle x^n\rangle\),
\begin{equation}
 n m_{n-1}-(n+3)m_{n+1}+\kappa(m_n-m_{n+2})=0,
 \quad 3m_1=\kappa(1-m_1^2-v),
 \quad v\ge\tfrac14e^{-2|\kappa|}>0.
 \qquad\text{\workstar}\label{early:closure}
\end{equation}
The last inequality follows from domination of Haar density by
\(e^{-2|\kappa|}\) and \(\operatorname{Var}(x)=\inf_c\langle(x-c)^2\rangle\).
Discarding \(v\) while keeping the true mean leaves residual \(\kappa v\).
At \(\kappa=0\), the next equation, not the first, enforces \(m_2=1/4\).
Adding a fluctuation symbol does not compute its next moment.
The same audit checked specific displayed formulas of
\texttt{arXiv:2608.05415v1} \cite{chatterjee2026lattice}: directed versus unordered plaquette counting,
a derivative prefactor, a missing source phase, and the rank impossibility
\(\eta^Tg\eta=I_{N_c^2-1}\) when \(N_c^2-1>D\).
These are local, version-specific tests and do not refute all closures.
Source: \repo{research/round12/closure/closure-audit.md}.

\paragraph{Change in understanding and the next obligations.}
By Round 12 the reliable object had narrowed from an intended coupled
Einstein--QED system to explicit finite models with named observables, domains,
and error norms.  The abelian branch still required a common renormalized
current and anisotropic stress for a specified state, followed by fluctuation
and regulator control.  The nonabelian branch had genuine infinite-representation
results on finite graphs, including interacting stationary and driven bounds.
Its remaining tasks were connected-volume estimates for the actual physical
Hamiltonian, scale matching, nontrivial continuum observables, reconstruction,
and a spectral exclusion covering the whole physical sector.  Neither the
formal proof planner nor the number of passed checks supplied those premises.
The machine-readable early ledger preserves every historical entry, retained
failure, source path, and contribution mapping; the present manuscript synthesis
does not claim to rerun every historical calculation.
